\chapter{From ideal access to a physical detector}
\label{ch:dg-entry}

The preceding geometry theorem answers a realization question: what must
every admissible hidden diffusion dissipate to produce a given ideal
observation law? A detector changes the law available to the scientist.
Exposure averages a signal over time, sampling retains selected moments,
and read noise adds uncertainty. Those operations must be included before
an ideal-law theorem becomes an experimental claim.

This extension separates three questions. Does the complete detector law
still identify the ideal law? Can a finite experiment estimate the desired
property with a stated error? How do preparation, read precision, duration
and measurement energy change that error? Different answers to these
questions can hold for the same physical model.

\section{One target, several kinds of access}

Use the stationary cosine-circle family as a reference:
\[
 d\phi_t=\omega\,dt+\sqrt{2\gamma}\,dW_t,
 \qquad Y_t=\cos\phi_t,\qquad \gamma>0.
\]
Its ideal scalar law is reversible. Within the ordinary hidden-diffusion
class of the preceding chapters, its exact compatible minimum is
$\mathcal M=\omega^2/\gamma$. This target is a property of the ideal
law, rather than an attempt to count arbitrary unobserved dissipative
subsystems. The detector considered here retains an exposure average plus
calibrated additive noise.

\begin{center}
\begin{tabular}{p{0.29\textwidth}p{0.60\textwidth}}
\toprule
Available access & Supported conclusion\\
\midrule
Complete densely sampled law; known exposure and independent noise &
The ideal law can be identified under the stated mixing assumptions.
The inverse is an infinite limiting construction.\\[3pt]
Fixed sampling grid after known exposure &
Distinct mixing diffusion mechanisms can have identical complete sampled
laws; an explicit Gaussian example proves this ambiguity.\\[3pt]
Finite stationary trials at fixed positive read noise &
No uniformly finite upper confidence bound exists for the circle minimum
over an unbounded frequency class.\\[3pt]
Known binary hypotheses with fixed read precision &
Passive stationary trials require and suffice at order $\omega^4$;
oriented phase preparation changes this order to $\omega^2$.\\[3pt]
Read precision purchased through the declared law $\nu^2=\kappa/E$ &
Optimized passive and prepared measurement-energy requirements both have
order $\omega^2$. Other apparatus costs remain additional.\\
\bottomrule
\end{tabular}
\end{center}

The count comparison concerns the known pair $\omega$ and $2\omega$
along a specified resonant sequence, at fixed positive diffusion, exposure,
read-noise level and target error. It does not solve estimation over an
arbitrary frequency interval. Preparation means setting the actual phase
to $\pi/2$; selecting a scalar value $Y=0$ does not resolve the two possible
orientations and does not supply the same operation.

\section{Why the resource distinction matters}

Stationary preparation makes the first moment of the exposure vanish.
The leading contrast through a noisy detector then appears in a second
moment. Oriented preparation supplies a first-moment contrast. This familiar
statistical mechanism explains the different trial-count exponents in the
proved example. It is not presented as a new general information law.

The same example shows why a smaller number of trials is insufficient to
establish an energy advantage. Under the explicitly declared sensor model,
passive experiments can spend more energy on read precision and use fewer
trials. The two access classes then share the same measurement-energy
exponent. The result compares orders, not equal optimal constants, and
does not include preparation, resetting, maintaining the drive, exposure
overhead or computation in a total-energy optimum.

The population proofs, finite-access theorems, exact countermodel and
deterministic design table form one account. The new table contains evaluated
analytic bounds, not additional simulated or measured experiments. The
earlier frozen simulation bank remains a separate numerical validation of
the finite-resolution certificate. Independent internal review checks both
layers; external peer review and prospective detector calibration remain
open.

The next practical dependency is a detector-specific joint calibration
guarantee for local volatility and the short-time statistics used in the
entropy certificate. Existing motion-blur correction and noisy-diffusion
inference provide essential baselines. Their assumptions must be carried
through the physical entropy bound, including selection and dependence
between calibration and test records.
